\chapter{Родитель статического
\texorpdfstring{$SU(2)$}{SU(2)}-пропагатора композитного синглета}
\label{ch:v10-discrete-resolution-su2-static-mediator}

Нулевая мода безмассового посредника не обязана разрушать физический
статический отклик. Разложим двухвершинный клеточный носитель на постоянную
gauge-моду и ортогональную физическую моду:
\begin{equation}
 P_0=\frac12\begin{pmatrix}1&1\\1&1\end{pmatrix},\qquad
 P_\perp=I-P_0,\qquad L=2P_\perp.
\end{equation}
При параметре gauge fixing $\xi>0$ оператор и его точный Green-оператор
равны
\begin{equation}
 K_\xi=2P_\perp+\xi P_0,
 \qquad G_\xi=\frac12P_\perp+\frac1\xi P_0,
 \qquad K_\xi G_\xi=I.
\end{equation}

Нормированный обменный источник
\begin{equation}
 j=\frac1{\sqrt2}(1,-1)^\top
\end{equation}
сохранён: $P_0j=0$ и $P_\perp j=j$. Поэтому физическая восприимчивость
\begin{equation}
 \chi_0=j^\top G_\xi j=\frac12,
 \qquad \frac{\partial\chi_0}{\partial\xi}=0
\end{equation}
конечна и не зависит от gauge fixing. Вместе с $g^2=3/8$ это условно даёт
\begin{equation}
 \kappa=g^2\chi_0=\frac3{16}.
\end{equation}

Родитель коэффициентов Green-оператора строг: для координат при
$(P_0,P_\perp)$ его гессиан равен $\operatorname{diag}(\xi^2,4)$ и имеет
определитель $4\xi^2>0$.

Но геометрический лапласиан ещё не является $SU(2)$-калибровочным полем.
В унаследованном носителе отсутствуют два морфизма: вложение adjoint-поля в
клеточные рёбра и отображение композитной пары в сохранённый ток $j$.
Поэтому числа $1/2$ и $3/16$ пока являются строгим условным следствием
архитектуры, а не физически выведенным связыванием.

\begin{theorem}[Сохранённый источник устраняет gauge-зависимость]
Проекция статического Green-оператора на ортогональный постоянной моде ток
равна $P_\perp/2$. Восприимчивость $1/2$ и коэффициент $3/16$ конечны и не
зависят от параметра gauge fixing.
\end{theorem}

\begin{nogocriterion}[Клеточный лапласиан не является gauge-полем сам по себе]
Совпадение операторной формы не доказывает, что рёберная переменная несёт
adjoint-представление $SU(2)$ и взаимодействует с дублетной парой через
сохранённый ток. Эти типовые морфизмы должны быть построены отдельно.
\end{nogocriterion}

\begin{protocol}\begin{enumerate}
\item условная архитектура: \textbf{$12/12$};
\item gauge-независимость: \textbf{$5/5$};
\item унаследованный геометрический лапласиан: \textbf{$1/1$};
\item типизированное $SU(2)$-вложение и ток: \textbf{$0/2$};
\item физическое происхождение: \textbf{$0/3$};
\item следующий гейт: \textbf{клеточное вложение $SU(2)$-посредника}.
\end{enumerate}\end{protocol}

Машинный сертификат сохранён в
\path{s2t_v10_cell_birth_four_volume_nonequilibrium_bath_discrete_resolution_transmuted_mode_asymptotically_free_su2_gauge_singlet_su2_mediator_static_propagator_parent_origin_gate_results.json}.